% Bagian: second-order-logic
% Bab: syntax-and-semantics
% Subbagian: satisfaction

\documentclass[../../../include/open-logic-section]{subfiles}

\begin{document}

\olfileid[id]{sol}{syn}{sat}

\olsection{Keterpenuhan}

\begin{explain}
Untuk mendefinisikan relasi keterpenuhan $\Sat{M}{!A}[s]$ bagi !!{formula}s
orde kedua, kita harus memperluas definisi-definisi tersebut agar mencakup
!!{variable}s orde kedua. Gagasan !!a{structure} sama bagi logika orde kedua
dan logika orde pertama. Perbedaannya hanya terletak pada penugasan
variabel~$s$: penugasan ini sekarang harus memberikan nilai bukan hanya bagi
!!{variable}s orde pertama, melainkan juga bagi !!{variable}s orde kedua.
\end{explain}

\begin{defn}[Penugasan Variabel]
Suatu \emph{penugasan variabel}~$s$ bagi !!a{structure}~$\Struct{M}$ ialah
fungsi yang memetakan setiap
\begin{enumerate}
\item !!{variable} objek~$\Obj{v_i}$ ke suatu anggota dari~$\Domain M$,
  yakni, $s(\Obj{v_i}) \in \Domain{M}$;
\item variabel relasi beraritas~$n$~$\Obj{V_i^n}$ ke suatu relasi
  beraritas~$n$ pada~$\Domain{M}$, yakni,
  $s(\Obj{V_i^n}) \subseteq \Domain{M}^n$;
\item variabel fungsi beraritas~$n$~$\Obj{u_i^n}$ ke suatu fungsi
  beraritas~$n$ dari $\Domain{M}^n$ ke~$\Domain{M}$, yakni,
  $s(\Obj{u_i^n})\colon \Domain{M}^n \to \Domain{M}$.
\end{enumerate}
\end{defn}

\begin{explain}
!!^a{structure} menetapkan !!a{value} kepada setiap !!{constant}, fungsi
kepada setiap !!{function}, dan relasi kepada setiap !!{predicate}; sedangkan
penugasan variabel orde kedua menetapkan objek, relasi, dan fungsi kepada
masing-masing variabel objek, relasi, dan fungsi. Bersama-sama, keduanya
memungkinkan kita menetapkan suatu nilai kepada setiap suku.
\end{explain}

\begin{defn}[\usetoken{S}{value} Suatu Suku]
Jika $t$ merupakan suku dari bahasa~$\Lang L$, $\Struct M$ merupakan
!!a{structure} bagi~$\Lang L$, dan $s$ merupakan penugasan
!!a{variable} bagi~$\Struct M$, maka \emph{!!{value}}~$\Value{t}{M}[s]$
didefinisikan seperti bagi suku-suku orde pertama, dengan tambahan klausa
berikut:
\begin{quote}
\indcase{t}{\Atom{u}{t_1, \ldots, t_n}}{
\[
\Value{\indfrm}{M}[s] = s(u)(\Value{t_1}{M}[s], \ldots,
\Value{t_n}{M}[s]).
\]}
\end{quote}
\end{defn}

\begin{defn}[Varian-$x$]
Jika $s$ merupakan penugasan !!a{variable} bagi !!a{structure}~$\Struct M$,
maka setiap penugasan !!{variable}~$s'$ bagi~$\Struct M$ yang hanya mungkin
berbeda dari~$s$ dalam nilai yang ditetapkannya kepada~$x$ disebut
\emph{varian-$x$} dari~$s$. Jika $s'$ merupakan varian-$x$ dari~$s$, kita
menulis $\varAssign{s'}{s}{x}$. (Demikian pula bagi variabel orde kedua $X$
atau~$u$.)
\end{defn}

\begin{defn}
  Jika $s$ merupakan penugasan !!a{variable} bagi !!a{structure}~$\Struct M$
  dan $m \in \Domain{M}$, maka penugasan~$\Subst{s}{m}{x}$ ialah penugasan
  !!{variable} yang didefinisikan oleh
  \[\Subst{s}{m}{x}(y) = \begin{cases}
    m & \text{jika } y \ident x\\
    s(y) & \text{jika tidak}.
  \end{cases}\]
  Jika $X$ merupakan !!{variable} relasi beraritas~$n$ dan $R \subseteq
  \Domain{M}^n$, maka $\Subst{s}{R}{X}$ ialah penugasan !!{variable} yang
  didefinisikan oleh
  \[\Subst{s}{R}{X}(y) = \begin{cases}
    R & \text{jika } y \ident X\\
    s(y) & \text{jika tidak}.
  \end{cases}\]
  Jika $u$ merupakan !!{variable} fungsi beraritas~$n$ dan
  $f\colon \Domain{M}^n \to \Domain{M}$, maka $\Subst{s}{f}{u}$ ialah
  penugasan !!{variable} yang didefinisikan oleh
  \[\Subst{s}{f}{u}(y) = \begin{cases}
    f & \text{jika } y \ident u\\
    s(y) & \text{jika tidak}.
  \end{cases}\]
  Dalam setiap kasus, $y$ dapat berupa sebarang !!{variable} orde pertama
  atau orde kedua.
\end{defn}

\begin{defn}[Keterpenuhan]
Bagi !!{formula}s orde kedua~$!A$, definisi keterpenuhan sama seperti
\olref[fol][syn][sat]{defn:satisfaction}, dengan penambahan berikut:
\begin{enumerate}
\item \indcase{!A}{\Atom{X}{t_1, \dots, t_n}}{$\Sat{M}{\indfrm}[s]$
  jika dan hanya jika $\langle \Value{t_1}{M}[s], \dots,
  \Value{t_n}{M}[s] \rangle \in s(X)$.}
\tagitem{prvAll}{%
  \indcase{!A}{\lforall[X][!B]}{$\Sat{M}{\indfrm}[s]$ jika dan hanya jika
  bagi setiap $R \subseteq \Domain{M}^n$,
  $\Sat{M}{!B}[\Subst{s}{R}{X}]$.}}{}

\tagitem{prvEx}{%
  \indcase{!A}{\lexists[X][!B]}{$\Sat{M}{\indfrm}[s]$ jika dan hanya jika
  bagi setidaknya satu $R \subseteq \Domain{M}^n$ berlaku
  $\Sat{M}{!B}[\Subst{s}{R}{X}]$.}}{}

\tagitem{prvAll}{%
  \indcase{!A}{\lforall[u][!B]}{$\Sat{M}{\indfrm}[s]$ jika dan hanya jika
  bagi setiap $f\colon \Domain{M}^n \to \Domain{M}$,
  $\Sat{M}{!B}[\Subst{s}{f}{u}]$.}}{}

\tagitem{prvEx}{%
  \indcase{!A}{\lexists[u][!B]}{$\Sat{M}{\indfrm}[s]$ jika dan hanya jika
  bagi setidaknya satu $f\colon \Domain{M}^n \to \Domain{M}$ berlaku
  $\Sat{M}{!B}[\Subst{s}{f}{u}]$.}}{}
\end{enumerate}
\end{defn}

\begin{ex}
  Perhatikan !!{formula} $\lforall[z][(\Atom{X}{z} \liff \lnot
    \Atom{Y}{z})]$. Formula ini tidak memuat kuantor orde kedua, tetapi
  memuat !!{variable}s orde kedua $X$ dan~$Y$ (di sini dipahami beraritas
  satu). !!{sentence} orde pertama yang bersesuaian,
  $\lforall[z][(\Atom{P}{z} \liff \lnot \Atom{R}{z})]$, menyatakan bahwa
  apa pun yang termasuk dalam interpretasi~$P$ tidak termasuk dalam
  interpretasi~$R$, dan sebaliknya. Dalam !!a{structure}, interpretasi
  !!a{predicate}~$P$ diberikan oleh interpretasi~$\Assign{P}{M}$. Namun,
  bagi !!{variable}s orde kedua seperti $X$ dan~$Y$, interpretasinya diberikan
  bukan oleh !!{structure} itu sendiri, melainkan oleh suatu penugasan
  !!a{variable}. Karena !!{formula} orde kedua tersebut bukan
  !!a{sentence} (formula itu memuat !!{variable}s bebas $X$ dan~$Y$), formula
  itu hanya terpenuhi relatif terhadap !!a{structure}~$\Struct{M}$ bersama
  suatu penugasan !!a{variable}~$s$.

  $\Sat{M}{\lforall[z][(\Atom{X}{z} \liff \lnot \Atom{Y}{z})]}[s]$
  berlaku tepat ketika !!{element}s dari~$s(X)$ bukan !!{element}s
  dari~$s(Y)$, dan sebaliknya, yakni, jika dan hanya jika
  $s(Y) = \Domain{M} \setminus s(X)$. Misalnya, ambil
  $\Domain{M} = \{1, 2, 3\}$. Karena tidak ada !!{predicate}s,
  !!{function}s, ataupun !!{constant}s yang terlibat, hanya domain
  dari~$\Struct{M}$ yang relevan. Sekarang, bagi $s_1(X) = \{1, 2\}$ dan
  $s_1(Y) = \{3\}$, kita mempunyai
  $\Sat{M}{\lforall[z][(\Atom{X}{z} \liff \lnot
      \Atom{Y}{z})]}[s_1]$.

  Sebaliknya, jika $s_2(X) = \{1, 2\}$ dan $s_2(Y) = \{2, 3\}$, maka
  $\Sat/{M}{\lforall[z][(\Atom{X}{z} \liff \lnot
      \Atom{Y}{z})]}[s_2]$. Alasannya,
  $\Sat{M}{\Atom{X}{z}}[\Subst{s_2}{2}{z}]$ (karena
  $2 \in \Subst{s_2}{2}{z}(X)$), tetapi
  $\Sat/{M}{\lnot \Atom{Y}{z}}[\Subst{s_2}{2}{z}]$ (karena juga
  $2 \in \Subst{s_2}{2}{z}(Y)$).
\end{ex}

\begin{ex}
  $\Sat{M}{\lexists[Y][(\lexists[y][\Atom{Y}{y}] \land
  \lforall[z][(\Atom{X}{z} \liff \lnot \Atom{Y}{z})])]}[s]$ berlaku jika
  terdapat $N \subseteq \Domain{M}$ sedemikian sehingga
  $\Sat{M}{(\lexists[y][\Atom{Y}{y}] \land \lforall[z][(\Atom{X}{z}
  \liff \lnot \Atom{Y}{z})])}[\Subst{s}{N}{Y}]$. Hal itu terjadi bagi
  sebarang $N \neq \emptyset$ (sehingga
  $\Sat{M}{\lexists[y][\Atom{Y}{y}]}[\Subst{s}{N}{Y}]$) yang, seperti dalam
  contoh sebelumnya, memenuhi $N = \Domain{M} \setminus s(X)$. Dengan kata
  lain,
  $\Sat{M}{\lexists[Y][(\lexists[y][\Atom{Y}{y}] \land
  \lforall[z][(\Atom{X}{z} \liff \lnot \Atom{Y}{z})])]}[s]$ jika dan hanya
  jika $\Domain{M} \setminus s(X)$ tidak kosong, yakni,
  $s(X) \neq \Domain{M}$. Jadi, !!{formula} tersebut terpenuhi, misalnya,
  jika $\Domain{M} = \{1, 2, 3\}$ dan $s(X) = \{1, 2\}$, tetapi tidak jika
  $s(X) = \{1, 2, 3\} = \Domain{M}$.

  Karena !!{formula} itu tidak terpenuhi setiap kali $s(X) = \Domain{M}$,
  !!{sentence}
  \[
  \lforall[X][\lexists[Y][(\lexists[y][\Atom{Y}{y}] \land
      \lforall[z][(\Atom{X}{z} \liff \lnot \Atom{Y}{z})])]]
  \]
  tidak pernah terpenuhi: bagi setiap !!{structure}~$\Struct{M}$, penugasan
  $s(X) = \Domain{M}$ akan membuat !!{sentence} tersebut salah. Di sisi lain,
  kalimat
  \[
  \lexists[X][\lexists[Y][(\lexists[y][\Atom{Y}{y}] \land
      \lforall[z][(\Atom{X}{z} \liff \lnot \Atom{Y}{z})])]]
  \]
  terpenuhi relatif terhadap setiap penugasan~$s$, karena kita selalu dapat
  menemukan $N \subseteq \Domain{M}$ dengan $N \neq \Domain{M}$ (misalnya,
  $N = \emptyset$).
\end{ex}

\begin{ex}
  !!{sentence} orde kedua~$\lforall[X][\lforall[y][X(y)]]$ menyatakan bahwa
  setiap relasi beraritas~$1$, yakni setiap sifat, berlaku atas setiap objek.
  Hal itu jelas tidak pernah benar, sebab dalam setiap~$\Struct{M}$, bagi
  penugasan variabel~$s$ dengan $s(X) = \emptyset$ dan
  $s(y) = a \in \Domain{M}$, kita mempunyai $\Sat/{M}{X(y)}[s]$. Ini berarti
  bahwa $!A \lif \lforall[X][\lforall[y][X(y)]]$ ekuivalen dalam logika orde
  kedua dengan $\lnot !A$, yakni: $\Sat{M}{!A \lif
  \lforall[X][\lforall[y][X(y)]]}$ jika dan hanya jika
  $\Sat{M}{\lnot !A}$. Dengan kata lain, dalam logika orde kedua kita dapat
  mendefinisikan $\lnot$ dengan menggunakan $\lforall$ dan~$\lif$.
\end{ex}

\begin{prob}
  Tunjukkan bahwa dalam logika orde kedua $\lforall$ dan~$\lif$ dapat
  mendefinisikan penghubung-penghubung lainnya:
  \begin{enumerate}
    \item Buktikan bahwa dalam logika orde kedua $!A \land !B$ ekuivalen
    dengan $\lforall[X][(!A \lif (!B \lif \lforall[x][X(x)])\lif
    \lforall[x][X(x)])]$.
    \item Temukan suatu formula orde kedua yang hanya menggunakan $\lforall$
    dan~$\lif$ serta ekuivalen dengan $!A \lor !B$.
  \end{enumerate}
\end{prob}

\end{document}
